% ============================================================
%  Chapter 5 / Section V
%  RTGS-Based Microservice Deployment
%
%  IEEE Transactions style  |  English  |  third-person passive
%
%  Notation (defined here; supersedes earlier sections where noted):
%    \mathcal{G}_t = (\mathcal{V}_t, \mathcal{E}_t)
%        Task graph  — subscript 't' denotes *task*, NOT time.
%        v_i \in \mathcal{V}_t : microservice (task) node
%        e_{i,j} \in \mathcal{E}_t : data-dependency edge
%
%    \mathcal{G}_r = (\mathcal{V}_r, \mathcal{E}_r)
%        UAV resource graph  — subscript 'r' denotes *resource*.
%        u_k \in \mathcal{V}_r : UAV node
%
%    n  — decision epoch (integer step index, avoids clash with
%          the task subscript 't')
%    N  = |\mathcal{V}_t|  — current task count (N \leq N_{\max}=20)
%    M  = |\mathcal{V}_r|  — fleet size (M = 8)
%
%  Required preamble packages:
%    \usepackage{amsmath, amssymb, bm}
%    \usepackage{algorithm, algpseudocode}
%    \usepackage{booktabs, multirow}
%    \usepackage{enumitem}
% ============================================================

% ----------------------------------------------------------------
\section{RTGS-Based Microservice Deployment}
\label{sec:microservice}

Edge-computing deployments on heterogeneous UAV swarms introduce a
class of workload management problems that differs fundamentally from
conventional cloud scheduling.
Microservice pipelines—containerised functional units connected by
explicit data-flow dependencies—map naturally onto DAGs, yet the
physical substrate on which they execute is radically non-stationary:
UAV bandwidth fluctuates with inter-vehicle distance, compute
capacity shrinks as batteries deplete, and the active set of
executable nodes contracts over time.
Static graph-to-UAV mapping strategies that treat the microservice
DAG as a fixed input fail to exploit this dynamism;
they cannot, for instance, consolidate two lightweight services to
eliminate a cross-UAV transfer when the intervening wireless link
degrades, nor decompose a bottleneck service to exploit a suddenly
available high-CPU UAV.

This section instantiates the RTGS framework introduced in
Section~\ref{sec:method} for the microservice deployment problem.
The task graph $\mathcal{G}_t = (\mathcal{V}_t, \mathcal{E}_t)$
represents the microservice pipeline, where each node
$v_i \in \mathcal{V}_t$ is one deployable service instance and each
directed edge $e_{i,j} \in \mathcal{E}_t$ encodes a data dependency
requiring intermediate results to be transferred from the UAV
executing $v_i$ to the UAV executing $v_j$ before $v_j$ may start.
The UAV swarm is modelled as a resource graph
$\mathcal{G}_r = (\mathcal{V}_r, \mathcal{E}_r)$, where
$\mathcal{V}_r = \{u_1, \ldots, u_M\}$ is the fleet and
$\mathcal{E}_r$ is the fully-connected edge set representing all
pairwise communication channels.

The remainder of this section is organised as follows.
Section~\ref{subsec:mdp} formalises the MDP over the joint state of
$\mathcal{G}_t$ and $\mathcal{G}_r$.
Section~\ref{subsec:gnn_attn} describes how task and resource
features are extracted and fused via dual-graph GNNs and a
cross-graph attention mechanism.
Section~\ref{subsec:operators} derives the graph-theoretic
properties of the Split and Merge shaping operators.
Section~\ref{subsec:ppo_mask} presents the PPO training procedure
with structural action masking that guarantees all graph
transformations preserve DAG acyclicity.

% ----------------------------------------------------------------
\subsection{MDP Formulation}
\label{subsec:mdp}

The RTGS microservice deployment problem is formalised as a finite-horizon
Markov Decision Process
$\mathcal{M} = \langle\mathcal{S},\,\mathcal{A},\,\mathcal{R},\,\gamma\rangle$.

% - - - - - - - - - - - - - - - - - - - - - - - - - - - - - - - -
\subsubsection{State Space $\mathcal{S}$}
\label{subsubsec:state}

A state $s \in \mathcal{S}$ at decision epoch $n$ is a composite
observation:
\begin{equation}
  s =
  \bigl(
    \mathbf{X}_t,\;\mathbf{E}_t,\;
    \mathbf{X}_r^{(n)},\;\mathbf{E}_r,\;
    \mathbf{f}^{(n)}
  \bigr),
  \label{eq:state_full}
\end{equation}
where $\mathbf{X}_t$ and $\mathbf{E}_t$ encode the current microservice
graph, $\mathbf{X}_r^{(n)}$ encodes the time-varying UAV resource state,
$\mathbf{E}_r$ is the fixed resource-graph edge index, and
$\mathbf{f}^{(n)}$ is the scheduling feedback vector from the preceding
epoch.
Each component is detailed below.

\paragraph{Task-Graph Features}

The microservice graph $\mathcal{G}_t$ is encoded as a zero-padded
node-feature matrix
$\mathbf{X}_t \in \mathbb{R}^{N_{\max} \times 2}$,
where $N_{\max} = 20$ is the maximum admissible service count.
The feature vector of service $v_i$ is
\begin{equation}
  \mathbf{x}_i^t = \bigl[c_i,\; d_i\bigr]^\top,
  \label{eq:task_feat}
\end{equation}
where $c_i > 0$ (CPU units) is the compute demand of $v_i$ and
$d_i > 0$ (MB) is its local data volume.
Each dependency edge $e_{i,j} \in \mathcal{E}_t$ additionally carries
a transfer volume $\delta_{ij} \geq 0$ (MB), representing the payload
that must be communicated from the executor of $v_i$ to the executor
of $v_j$.
The edge set is stored as a COO-format index
$\mathbf{E}_t \in \mathbb{Z}^{2 \times N_{\max}^2}$, zero-padded
to a fixed buffer.
Padding rows in $\mathbf{X}_t$ are identically zero and are excluded
from all graph operations via Boolean masking.

\paragraph{Resource-Graph Features}

Unlike the task graph, the resource-graph feature matrix
$\mathbf{X}_r^{(n)} \in [0,1]^{M \times 3}$ is \emph{re-computed at
every epoch} by the mobility model, reflecting the current physical
state of the swarm.
The feature vector of UAV $u_k$ at epoch $n$ is
\begin{equation}
  \mathbf{x}_k^r(n) =
  \Bigl[
    \underbrace{\tilde{\phi}_k(n)\,/\,\phi^{\max}}_{\text{CPU ratio}},\quad
    \underbrace{\bar{B}_k(n)\,/\,B_{\max}}_{\text{BW ratio}},\quad
    \underbrace{\beta_k(n)}_{\text{battery}}
  \Bigr]^\top \!\in [0,1]^3,
  \label{eq:res_feat}
\end{equation}
where the three components are defined as follows.

\textit{Battery-derated CPU ratio.}
The effective CPU capacity of $u_k$ is a piecewise function of its
residual battery $\beta_k(n) \in [0,1]$:
\begin{equation}
  \tilde{\phi}_k(n) =
  \begin{cases}
    0,          & \beta_k(n) = 0
                  \;\text{(depleted — offline)},\\
    \gamma\,\phi_k, & 0 < \beta_k(n) < \beta_{\mathrm{low}}
                  \;\text{(low-power mode, \(\gamma = 0.5\))},\\
    \phi_k,     & \beta_k(n) \geq \beta_{\mathrm{low}}
                  \;\text{(nominal)},
  \end{cases}
  \label{eq:cpu_derate}
\end{equation}
with low-power threshold $\beta_{\mathrm{low}} = 0.20$ and nominal
CPU capacity $\phi_k$ drawn from the heterogeneous fleet profile.
The CPU ratio is then $\tilde{\phi}_k(n)/\phi^{\max}$, where
$\phi^{\max} = \max_k \phi_k$ is the fleet-wide ceiling.

\textit{Mean link-bandwidth ratio.}
The pairwise link bandwidth between active UAVs $u_k$ and $u_l$ is
governed by an exponential distance-decay model:
\begin{equation}
  B_{kl}(n) =
  \begin{cases}
    B_{\min} + \bigl(B_{\max} - B_{\min}\bigr)\,
      e^{-r_{kl}(n)/\delta_H},
      & r_{kl}(n) \leq D_c, \\
    B_{\min},
      & r_{kl}(n) > D_c,
  \end{cases}
  \label{eq:bw_model}
\end{equation}
where $r_{kl}(n)$ is the Euclidean inter-UAV distance at epoch $n$,
$\delta_H = 150\,\text{m}$ is the characteristic decay distance,
$D_c = 400\,\text{m}$ is the communication radius,
$B_{\max} = 1600\,\text{Mbps}$, and $B_{\min} = 200\,\text{Mbps}$.
The mean-bandwidth feature is
\begin{equation}
  \bar{B}_k(n) =
  \frac{1}{|\mathcal{U}_a(n)| - 1}
  \sum_{\substack{u_l \in \mathcal{U}_a(n) \\ l \neq k}}
  B_{kl}(n),
  \label{eq:avg_bw}
\end{equation}
where $\mathcal{U}_a(n) = \{u_k : \beta_k(n) > 0\}$ is the active
UAV set at epoch $n$.
All three feature dimensions are normalised to $[0,1]$,
yielding a scale-invariant input for the GNN encoder.

\paragraph{Scheduling Feedback Vector}

The six-dimensional feedback vector
$\mathbf{f}^{(n)} \in \mathbb{R}^6$ summarises the outcome of the
greedy scheduler applied to the task graph at the end of epoch $n$:
\begin{equation}
  \mathbf{f}^{(n)} =
  \Bigl[
    \underbrace{\lambda^{(n)}}_{\text{latency ratio}},\;
    \underbrace{\sigma^{(n)}}_{\text{success rate}},\;
    \underbrace{\varrho^{(n)}}_{\text{reschedule count}},\;
    \underbrace{\bar{u}^{(n)}}_{\text{fleet util.}},\;
    \underbrace{\hat{B}_{\min}^{(n)}}_{\text{min BW ratio}},\;
    \underbrace{\hat{\beta}_{\min}^{(n)}}_{\text{min battery}}
  \Bigr]^\top,
  \label{eq:feedback_vec}
\end{equation}
where:
\begin{itemize}[leftmargin=*, itemsep=1pt]
  \item $\lambda^{(n)} = T_{\mathrm{span}}^{(n)}\,/\,T_D \in \mathbb{R}_{\geq 0}$
        is the schedule makespan normalised by the deadline $T_D$;
        $\lambda > 1$ indicates a deadline miss.
  \item $\sigma^{(n)} \in [0,1]$
        is the fraction of services successfully assigned to an
        active UAV with sufficient residual CPU.
  \item $\varrho^{(n)} \in \mathbb{Z}_{\geq 0}$
        counts services that required a forced assignment to an
        overloaded UAV (partial execution).
  \item $\bar{u}^{(n)} \in [0,1]$
        is the mean CPU utilisation across all active UAVs.
  \item $\hat{B}_{\min}^{(n)} = \min_{k \neq l,\,
        u_k,u_l \in \mathcal{U}_a} B_{kl}(n)\,/\,B_{\max} \in [0,1]$
        is the weakest active inter-UAV link, providing a
        mobility-induced fragmentation signal.
  \item $\hat{\beta}_{\min}^{(n)} = \min_k \beta_k(n) \in [0,1]$
        is the most-depleted UAV's battery fraction,
        signalling impending loss of fleet nodes.
\end{itemize}
The joint inclusion of link-quality and battery signals in the state
is a deliberate design choice: both are non-stationary quantities
that directly determine whether Split (exploiting available CPU) or
Merge (eliminating weak-link transfers) is the more profitable action
at epoch $n$.

% - - - - - - - - - - - - - - - - - - - - - - - - - - - - - - - -
\subsubsection{Action Space $\mathcal{A}$}
\label{subsubsec:action}

The action set $\mathcal{A}$ is a flat discrete space of
\emph{node-level graph shaping operations}:
\begin{equation}
  \mathcal{A} =
  \underbrace{\{a_{\varnothing}\}}_{\text{noop}}
  \;\cup\;
  \underbrace{\bigl\{a_i^{\mathrm{sp}} : i \in [N_{\max}]\bigr\}}_{\text{Split actions}}
  \;\cup\;
  \underbrace{\bigl\{a_{ij}^{\mathrm{me}} :
    0 \leq i < j < N_{\max}\bigr\}}_{\text{Merge actions}},
  \label{eq:action_set}
\end{equation}
with total cardinality
\begin{equation}
  |\mathcal{A}| = 1 + N_{\max} + \binom{N_{\max}}{2}
               = 1 + 20 + 190 = 211.
  \label{eq:action_dim}
\end{equation}

In practice, actions are represented as a single integer
$a \in \{0, 1, \ldots, 210\}$ and decoded via a fixed lookup table
into a structured triple $(op,\,p_i,\,p_j)$, where $op$ identifies
the operation type and $p_i,\,p_j \in [N_{\max}]$ are
\emph{padded node positions} (indices into the sorted node list of
the current graph):
\begin{equation}
  \mathrm{decode}(a) =
  \begin{cases}
    (\mathrm{noop},\; -,\; -)
      & a = 0, \\[3pt]
    (\mathrm{split},\; a - 1,\; -)
      & 1 \leq a \leq N_{\max}, \\[3pt]
    \bigl(\mathrm{merge},\; p_i^{(a)},\; p_j^{(a)}\bigr)
      & N_{\max} < a \leq 210,
  \end{cases}
  \label{eq:decode}
\end{equation}
where $(p_i^{(a)},\, p_j^{(a)})$ is the $(a - N_{\max} - 1)$-th
entry of the pre-enumerated lexicographic list of all
$\binom{20}{2} = 190$ ordered pairs $(i,j)$, $0 \leq i < j < 20$.

\begin{table}[t]
  \centering
  \caption{Decomposition of the 211-Dimensional Action Space}
  \label{tab:action_space}
  \begin{tabular}{@{}lccl@{}}
    \toprule
    Type & Index range & Count & Semantics \\
    \midrule
    Noop   & $0$              & $1$   & No topology change \\
    Split  & $1$ -- $20$      & $20$  & Split node at padded position $a-1$ \\
    Merge  & $21$ -- $210$    & $190$ & Merge pair $(p_i^{(a)}, p_j^{(a)})$,\;
                                        $\binom{20}{2}$ lex.\ pairs \\
    \midrule
    Total  & $0$ -- $210$     & $\mathbf{211}$ & \\
    \bottomrule
  \end{tabular}
\end{table}

Table~\ref{tab:action_space} summarises the index layout.
Node-level granularity is the key distinction from prior work
that restricts the shaping policy to a coarse three-way choice
(noop / split-heaviest / merge-cheapest): here the agent
selects \emph{which specific service} to restructure, conditioned
on cross-graph attention weights that reflect the current UAV
resource context (see Section~\ref{subsec:gnn_attn}).
This fine-grained design expands the action space by two orders of
magnitude relative to a coarse binary policy, but the increase is
managed through structural action masking
(Section~\ref{subsec:ppo_mask}) that eliminates infeasible actions
before sampling.

% - - - - - - - - - - - - - - - - - - - - - - - - - - - - - - - -
\subsubsection{Reward Function $\mathcal{R}$}
\label{subsubsec:reward}

The reward $r^{(n)}$ at epoch $n$ is composed of three semantically
distinct terms:
\begin{equation}
  \boxed{
  r^{(n)} =
    \underbrace{r_{\mathrm{qual}}^{(n)}}_{\text{absolute quality}}
    +\;
    \underbrace{r_{\Delta}^{(n)}}_{\text{incremental gain}}
    +\;
    \underbrace{r_{\mathrm{pen}}^{(n)}}_{\text{penalty}}
  }
  \label{eq:reward_decomp}
\end{equation}

\paragraph{Absolute Quality Term $r_{\mathrm{qual}}$}

This term rewards the agent for achieving high service success rates
and feasible latencies at every epoch, regardless of the previous
state.
It ensures that a policy which \emph{maintains} excellent performance
continues to receive positive reinforcement:
\begin{equation}
  r_{\mathrm{qual}}^{(n)} =
    \underbrace{2\,\sigma^{(n)}}_{\text{success rate}}
    \;+\;
    \underbrace{1.5\,\max\!\bigl(0,\;1 - \lambda^{(n)}\bigr)}_{\text{latency margin}}.
  \label{eq:r_qual}
\end{equation}
The coefficient $2$ on the success-rate term is larger than the
$1.5$ latency coefficient, reflecting the priority that all services
must be deployable over merely meeting the deadline.

\paragraph{Incremental Improvement Term $r_{\Delta}$}

This term provides dense shaping by rewarding step-to-step
improvements in both scheduling dimensions, preventing the policy
from stagnating in any mediocre-but-stable local optimum:
\begin{equation}
  r_{\Delta}^{(n)} =
    \underbrace{
      \bigl(\lambda^{(n-1)} - \lambda^{(n)}\bigr)
    }_{\text{latency reduction}}
    +\;
    \underbrace{
      0.5\,\bigl(\sigma^{(n)} - \sigma^{(n-1)}\bigr)
    }_{\text{success gain}},
  \label{eq:r_delta}
\end{equation}
where $\lambda^{(n-1)}$ and $\sigma^{(n-1)}$ are the values from
the preceding epoch, initialised at episode start to $\lambda^{(0)}=1$
(at-deadline) and $\sigma^{(0)}=0$ (no success).
Both delta terms admit negative values, penalising performance
regression.
The asymmetric weight $(1.0\;\text{vs.}\;0.5)$ assigns greater
incentive to latency reduction, consistent with the latency-sensitive
nature of microservice pipelines.

\paragraph{Penalty Term $r_{\mathrm{pen}}$}

Three negative sub-terms discourage distinct failure modes:
\begin{equation}
  r_{\mathrm{pen}}^{(n)} =
    \underbrace{
      -\,\frac{\varrho^{(n)}}{N^{(n)}}
    }_{\substack{\text{reschedule} \\ \text{rate}}}
    \;-\;
    \underbrace{
      2\,\max\!\bigl(0,\;\lambda^{(n)} - 1\bigr)
    }_{\substack{\text{deadline} \\ \text{overshoot}}}
    \;-\;
    \underbrace{
      0.05\,\max\!\bigl(0,\;N^{(n)} - (N_{\max}-2)\bigr)
    }_{\substack{\text{topology} \\ \text{inflation}}},
  \label{eq:r_pen}
\end{equation}
where $N^{(n)} = |\mathcal{V}_t^{(n)}|$ is the task-graph size at
epoch $n$.

\textit{Reschedule-rate penalty.}
Rather than penalising the raw count $\varrho^{(n)}$, the term is
normalised by the current graph size $N^{(n)}$.
This removes a systematic bias: a graph with twice as many nodes
naturally produces twice as many absolute reschedule events even if
the per-node rescheduling rate is identical, so raw-count penalties
would unfairly discourage productive Split operations that increase
$N^{(n)}$.

\textit{Deadline-overshoot penalty.}
The penalty grows linearly with the fractional latency excess
$(\lambda^{(n)} - 1)$ rather than applying a hard threshold,
preserving gradient information throughout the overconstrained
regime.
The coefficient $2.0$ is calibrated to dominate the maximum
achievable absolute-quality term $r_{\mathrm{qual}}$, ensuring
that deadline satisfaction is the primary objective.

\textit{Topology-inflation penalty.}
A weak coefficient of $0.05$ activates only when
$N^{(n)} > N_{\max} - 2 = 18$, i.e., within two nodes of the hard
capacity ceiling.
This prevents degenerate schedules in which the agent continuously
applies Split to bloat the graph—a failure mode that can arise
when the split logit receives strong gradient without a size
constraint.
The threshold is set two below $N_{\max}$ rather than at $N_{\max}$
itself to preserve a safety buffer: the policy can still execute one
more Split operation before the hard constraint blocks it.

\paragraph{Combined Reward Expression}

Substituting \eqref{eq:r_qual}--\eqref{eq:r_pen} into
\eqref{eq:reward_decomp} gives the full scalar reward:
\begin{align}
  r^{(n)} =\;&
    2\,\sigma^{(n)}
    + 1.5\,\max\!\bigl(0,\;1-\lambda^{(n)}\bigr)
    \nonumber\\
    &+ \bigl(\lambda^{(n-1)} - \lambda^{(n)}\bigr)
    + 0.5\,\bigl(\sigma^{(n)} - \sigma^{(n-1)}\bigr)
    \nonumber\\
    &- \frac{\varrho^{(n)}}{N^{(n)}}
    - 2\,\max\!\bigl(0,\;\lambda^{(n)}-1\bigr)
    \nonumber\\
    &- 0.05\,\max\!\bigl(0,\;N^{(n)}-18\bigr).
  \label{eq:reward_full}
\end{align}
The discount factor $\gamma = 0.99$ is applied in computing the
return $R^{(n)} = \sum_{k \geq 0} \gamma^k r^{(n+k)}$.

\paragraph{Termination Conditions}

An episode terminates under two conditions.
\textit{Catastrophic failure}: the schedule makespan exceeds twice
the deadline, $T_{\mathrm{span}}^{(n)} > 2\,T_D$, indicating an
irrecoverable scheduling collapse (e.g., total fleet depletion).
\textit{Horizon truncation}: the step count reaches the maximum
episode length $N_{\mathrm{ep}} = 50$, at which point the episode
is truncated and a bootstrap value estimate replaces the terminal
return.
